%-------------------------------------------------------------------------------
%	WORK EXPERIENCE
%-------------------------------------------------------------------------------
\cvsection{Work Experience}

\begin{cventries}
  \cventry
    {Applied Scientist Intern}
    {Zillow Group}
    {Remote, USA}
    {May 2025 -- Aug. 2025}
    {
      \begin{cvitems}
        \item Built a home-purchase assistant: a multi-turn chatbot that answers buyer questions, tracks purchase intent over the conversation, and narrows a large listing inventory to matching properties.
        \item Post-trained the dialogue model on long multi-turn conversations, first with DPO on preference pairs over whole dialogues, then with GRPO using a trajectory-level reward rather than a per-turn reward.
        \item Built the reward model and the evaluation pipeline for this training, including rollout generation and conversation-level scoring for grounded property-search dialogue.
      \end{cvitems}
    }

  \cventry
    {Machine Learning Research Engineer Intern}
    {ScaleAI, Inc.}
    {New York City, NY, USA}
    {Jan. 2025 -- May 2025}
    {
      \begin{cvitems}
        \item Built precedent retrieval as a hybrid RAG system over constitutional bylaws and prior Qatari Supreme Court cassation rulings: dense embedding retrieval fused with sparse lexical matching, followed by reranking of the merged candidate set.
        \item Fine-tuned an Arabic LLM with SFT on data collected specifically for Qatari court cases, for event extraction and case summarization.
        \item Designed the Arabic information-extraction pipeline: OCR over scanned filings, then event extraction that links each event to its supporting evidence spans, feeding a structured case summarizer.
        \item Deployed the above as API services behind a judge-facing case-analysis dashboard for the Qatari judicial department.
      \end{cvitems}
    }

  \cventry
    {AI/ML Research Intern}
    {Apple}
    {Seattle, WA, USA}
    {May 2022 -- Aug. 2022}
    {
      \begin{cvitems}
        \item Built a conditional synthetic dialogue generator for Siri training data, targeting scenarios with little or no real usage data. Mentored by Dr.\ Alex Churchill.
        \item Simulated the user side of each conversation so the pipeline produced complete multi-turn dialogues rather than isolated utterances.
        \item Conditioned generation on intent and user behavior, allowing a target dataset to be specified by requirements, and ran it zero-shot without seed dialogues.
      \end{cvitems}
    }

  \cventry
    {Research Assistant}
    {Qatar Computing Research Institute}
    {Doha, Qatar}
    {Jul. 2019 -- Aug. 2021}
    {
      \begin{cvitems}
        \item Research assistant advised by Dr.\ Preslav Nakov and Prof.\ Giovanni Da San Martino.
        \item Introduced the task of detecting previously fact-checked claims, with a retrieve-then-rerank pipeline: BM25 candidate retrieval, Sentence-BERT semantic similarity, and a learning-to-rank reranker (ACL 2020). Extended it to document-level claim retrieval and to context modeling with Transformer encoders over surrounding tweets and co-mentions (Findings of EMNLP 2022, Findings of NAACL 2022, IJCAI 2021).
        \item Built \textit{Prta}, a public system for fine-grained propaganda detection in news articles, combining span-level sequence tagging and sentence-level classification with a BERT-based multi-granularity network (ACL 2020, Best Demo Award Honorable Mention). Extended it to memes as multimodal multi-label classification with text--image fusion models (ACL 2021, SemEval-2021).
        \item Built multilingual, multi-label classifiers on fine-tuned Transformer encoders for the factuality, harmfulness, and framing of COVID-19 claims across Arabic, English, and Bulgarian social media (ICWSM 2021, Findings of EMNLP 2021, RANLP 2021, COLING 2022).
        \item Built cross-lingual emotion classifiers for low-resource languages using multilingual encoders and machine translation for transfer (LREC 2022).
        \item Co-organized shared tasks on check-worthiness estimation and verified-claim retrieval at CLEF--CheckThat! 2020 and 2021, on persuasion techniques in memes at SemEval-2021 Task 6, and at NLP4IF-2021, releasing the datasets and evaluation setups.
      \end{cvitems}
    }

  % \cventry
  %   {Research Intern}
  %   {The Robotics Institute, Carnegie Mellon University}
  %   {Pittsburgh, PA, USA}
  %   {May 2018 -- Aug. 2018}
  %   {
  %     \begin{cvitems}
  %       \item Advised by \href{https://dblp.org/pid/53/204.html}{Dr. Christoph Mertz}.
  %       \item Worked on road-image segmentation for lane detection and identification.
  %       \item Explored deep learning and classical computer-vision methods to identify and segment road markings for autonomous vehicles.
  %     \end{cvitems}
  %   }

  % \cventry
  %   {Part-Time Research Assistant}
  %   {Carnegie Mellon University}
  %   {Pittsburgh, PA, USA}
  %   {May 2017 -- Jun. 2018}
  %   {
  %     \begin{cvitems}
  %       \item Advised by \href{https://scholar.google.com/citations?user=IWcGY98AAAAJ&hl=en}{Prof. Bhiksha Raj}.
  %       \item Developed two algorithms competitive with state-of-the-art for estimating classifier accuracy under limited labeling resources (active testing).
  %       \item Provided mathematical analysis and empirical validation; published at ICMLA 2018.
  %     \end{cvitems}
  %   }
\end{cventries}
